\section{Introduction}
\label{sec:intro}

Knowledge editing promises surgical updates to a language model's factual
associations without retraining. Locate-then-edit methods
\citep{meng2022rome,meng2023memit} realize this with a rank-one or low-rank
write to a single MLP down-projection at a ``factual-storage'' layer. The
write is effective at installing the target fact, but it is not local: it
also shifts the model's read-out on \emph{unrelated} facts, a phenomenon we
call collateral damage. Understanding, predicting, and removing that damage is
now the central reliability question for the method family.

A natural question is whether collateral damage is predictable
\emph{a priori} -- before running the edit -- from geometry alone. The
concurrent work CLaRE-ty \citep{clarety2026} answers affirmatively with an
\emph{empirical} single-layer forward-activation cosine, evaluated over
\clareNfacts{} facts, \clareNeditors{} locate-then-edit editors, and
\clareNmodels{} attention models, and -- by its authors' own statement --
correlational, with no causal mechanism. We take a different and complementary
angle: not a new predictor, but an \emph{account} of the mechanism, its
failure modes, and a causal test.

\paragraph{Scope (stated up front).} The \emph{signed} key-cosine$\rightarrow$damage
law we characterize is \textbf{Llama-family-specific}. We frame the paper as a
Llama-family mechanistic account \emph{plus} a cross-architecture dissociation
arm, not a universal predictor. Two laws with different reach recur throughout:
the \emph{signed} law (sign and rank of the key-cosine coupling) holds on the
Llama family and off it either goes null (gemma, Phi) or inverts (Qwen); the
\emph{magnitude} law ($|C|\rightarrow|\text{damage}|$) transfers on four of five
non-Llama families, with gemma the documented exception.

\paragraph{Contributions.}
\begin{itemize}
\item \textbf{A closed-form \SxC{} decomposition (\S\ref{sec:method}).} ROME's
  collateral perturbation to a probe factorizes as
  $S\cdot\|k_p\|\cdot|C|$: ``cosine ranks damage'' is built into the rank-one
  algebra. The product \SxC{} is the exact closed-form reduction of
  first-order gradient influence, obtainable with no backprop -- it explains
  \emph{why} cosine works and predicts \emph{where} it stops.
\item \textbf{A depth/regime transition (\S\ref{sec:regime}).} On the
  confound-clean within-probe metric, key-cosine beats matrix-norm-growth at
  L8/L10/L12 and norm-growth overtakes at L14; the coupling's sign tracks the
  sign of the damage regime across scale.
\item \textbf{An editor/architecture dissociation (\S\ref{sec:dissociation}).}
  The law is locate-then-edit-mechanism-specific: it dies for full-rank
  fine-tuning and for MEMIT, and is near-null/inverted off Llama -- each
  breakdown predicted by the same algebra.
\item \textbf{A causal test (\S\ref{sec:causal}).} AlphaEdit's null-space
  projection removes the geometry-predicted damage in proportion to key-cosine
  at every layer, and erases the Qwen inversion -- the causal pathway CLaRE
  explicitly leaves open. A held-out projector retires the circularity objection.
\item \textbf{Extension to a new edit type (\S\ref{sec:deletion}).} The same
  mechanism governs collateral damage from refusal-\emph{deletion} edits, and
  transfers to zsRE.
\end{itemize}

Code, per-cell result artifacts, and figure-generation scripts sufficient to
reproduce every number will be released upon publication.
